% Meta: 400–550 palabras
\section{Introducción}

El sector agroalimentario enfrenta una mayor demanda de alimentos, el cambio climático y brechas de productividad e inclusión financiera. Los Fideicomisos Instituidos en Relación con la Agricultura (FIRA) atienden ese contexto con apoyos a la innovación tecnológica. Dentro del subcomponente Emprendedores Innovadores, el Reto AgroCebada 2026 pide un modelo analítico del rendimiento de parcelas de cebada en Hidalgo, Puebla y Tlaxcala, y una herramienta para consultar el resultado \citep{fira2026}.

La percepción remota describe el cultivo a lo largo del ciclo mediante índices de vegetación, y los rásteres climáticos aportan lluvia y temperatura. Esas fuentes ya existen. El problema es convertirlas en un estimado por parcela que resista un conjunto de evaluación cuyas etiquetas permanecen ocultas. El dataset oficial reúne 197 polígonos del ciclo abril–octubre de 2025. El rendimiento, en toneladas por hectárea, se observa en 138 parcelas. Las 59 restantes son el 30\,\% que el comité calificará. El equipo no usó esas etiquetas: entrenó solo con las 138 observadas y produjo un estimado para las 197.

Este reporte documenta el modelo de entrega, llamado mezcla, y el dashboard que lo consulta. La mezcla promedia dos predicciones. La primera es un ensamble de LightGBM, bosque aleatorio y red elástica \citep{ke2017,breiman2001,zou2005} entrenado con 121 variables de satélite, clima, topografía y ubicación. La segunda es un LightGBM que solo usa la longitud. Cada miembro del ensamble pesa un sexto y el modelo de longitud pesa un medio, de modo que la predicción final es el promedio simple del ensamble y del modelo de longitud.

La validación separa dos preguntas que no deben confundirse. La validación cruzada estratificada por municipio estima el error allí donde también están las parcelas ocultas: es la cifra que representa lo que se califica. Dejar un municipio completo fuera, esquema que aquí se llama LOMO, mide el deterioro cuando la zona es nueva dentro de la misma región. La distinción importa porque, en esta caja de unos 45 km, el rendimiento cambia de oeste a este y una parte de la señal es geográfica.

El escrito sigue la estructura pedida por las bases. Después de la justificación y los objetivos se describe la metodología, se reportan los resultados y se responden las tres preguntas guía. Un apartado vincula el estimado con la inclusión financiera del sector agropecuario. Cierran las conclusiones, la bibliografía y dos anexos: el trabajo del equipo y la declaración de uso de inteligencia artificial. Todas las cifras de desempeño provienen de los artefactos congelados que carga la API, no de una corrida paralela. El código, las instrucciones de ejecución y el enlace de la aplicación acompañan este documento para que el jurado pueda repetir el cálculo y consultar las parcelas.
